\section{$\lambda/4$-Plate}

    \subsection{Theory of a $\lambda/4$-plate}
 A $\lambda/4$-plate, also known as a quarter-wave plate, is an optical device
that introduces a phase shift of $\pi/2$ (90 degrees) between the orthogonal
components of polarized light.

This phase shift arises from birefringence: the plate is made of an
anisotropic crystal, in which light polarized parallel to the optic
axis (extraordinary ray) experiences a different refractive index
$n_e$ than light polarized perpendicular to the optic axis (the
ordinary ray), which experiences $n_o$. Since the two polarization
components travel through the same physical thickness $d$ but at different
phase velocities $v_o = c/n_o$ and $v_e = c/n_e$, they accumulate different
optical path lengths, and hence a relative phase difference, when passing
through the crystal.

For a plane wave of vacuum wavelength $\lambda$ traveling a distance $d$
through the crystal, each component picks up a phase
\begin{equation}
    \phi_o = \frac{2\pi}{\lambda} n_o d, \qquad
    \phi_e = \frac{2\pi}{\lambda} n_e d.
\end{equation}
\noindent
The phase difference between the two components after traversing the plate
is therefore
\begin{equation}
    \delta = \phi_e - \phi_o
    = \frac{2\pi}{\lambda}\left(n_e - n_o\right) d
    = \frac{2\pi}{\lambda}\,\Delta n \, d,
\end{equation}
where $\Delta n = n_e - n_o$ is the birefringence of the material.\\
\noindent
For a quarter-wave plate, we require $\delta = \pi/2$. Setting this equal to the expression above and solving for
$d$:
\begin{equation}
    \frac{2\pi}{\lambda}\,\Delta n\, d = \frac{\pi}{2}
    \quad\Longrightarrow\quad
    d = \frac{\lambda}{4\,|\Delta n|}.
\end{equation}
\noindent
This is the minimum thickness. In practice, plates are often
made thicker for mechanical robustness
which still produces a phase difference of $\pi/2$ modulo $2\pi$, since
adding any integer number of full wavelengths of relative delay does not
change the resulting polarization state.

    \subsection{Experimental setup for a $\lambda/4$-plate}
    In this experiment a $\lambda/4$-plate made of quartz was used. The light of an incandescent lamp was first linearly polarized using a polarizer.
    The linearly polarized light then passed through the $\lambda/4$-plate which was oriented at the angles $\alpha = 0^\circ, -45^\circ, -90^\circ$ 
    with respect to the polarization direction of the incident light. After passing through the $\lambda/4$-plate, the light was analyzed using a
    analyzer which was rotated to observe the changes in intensity and polarization state of the transmitted light. \\
    Afterwards the $\lambda/4$-plate was oriented at $\alpha = -45^\circ$ and another identical $\lambda/4$-plate was placed in the ray path, such that
    its optical axis was oriented parallel to the first $\lambda/4$-plate.

    \subsection{Results of a $\lambda/4$-plate}

    For $\alpha = 0^\circ$ and $\alpha = -90^\circ$, the transmitted light remained linearly polarized because the light is polarized along the same axis to
    one of the oplical axes and excites only that component. Thus no relative phase shift occurs and the polarization state does not change.
    For $\alpha = -45^\circ$, the transmitted light became circularly polarized because the light is polarized at 45 degrees to the optical axes of
    the $\lambda/4$-plate, exciting both components equally. The $\lambda/4$-plate introduces a phase shift of $\pi/2$ between the two components, resulting
    in circular polarization. At any other angle $\alpha  = ]0, -90^\circ[ \quad \textbackslash -45^\circ$ the transmitted light became elliptically polarized
    because the light is polarized at an angle that is not aligned with the optical axes of the $\lambda/4$-plate, exciting both components but with different
    amplitudes. \\
    When the second $\lambda/4$-plate was added with its optical axis parallel to the first one, tthe system is equivalent to a $\lambda/2$-plate,
    which introduces a phase shift of $\pi$ between the two components. The rotation angle is twice the angle between the incident polarization and the optical
    axis of the plate. The light remains linearly polarized but with a rotated polarization plane. \\
    If the quartz $\lambda/4$-plate is replaced by a plate made of calcite, the thickness of the plate would need to be adjusted to achieve the same phase
    shift of $\pi/2$ because the birefringence of calcite is different from that of quartz.